\section*{Capítulo \ref{ch:g6:rocks-and-minerals} --- Rochas e minerais}

\begin{solution}{exo:g6:rocks-and-minerals:1}
Granito: ígnea. Arenito: sedimentar. Mármore: metamórfica. Basalto: ígnea.
Calcário: sedimentar. Ardósia: metamórfica.
\end{solution}

\begin{solution}{exo:g6:rocks-and-minerals:2}
O quartzo é um \omterm{def:g6:rocks-and-minerals:mineral}{mineral}; o granito é uma \omterm{def:g3:stones-pebbles-sand:rock}{rocha}. Um \omterm{def:g6:rocks-and-minerals:mineral}{mineral} é uma única
substância natural com os seus átomos arrumados de forma regular; uma \omterm{def:g3:stones-pebbles-sand:rock}{rocha}
é feita de um ou de vários \omterm{def:g6:rocks-and-minerals:mineral}{minerais} grudados uns nos outros (o granito:
feldspato, quartzo e mica).
\end{solution}

\begin{solution}{exo:g6:rocks-and-minerals:3}
Talco (1), calcita (3), quartzo (7), diamante (10). O diamante, o mais
duro, risca todos os outros.
\end{solution}

\begin{solution}{exo:g6:rocks-and-minerals:4}
``A \omterm{def:g3:stones-pebbles-sand:rock}{rocha} derretida vira uma \omterm{def:g6:rocks-and-minerals:rock-families}{rocha ígnea} por resfriamento. Um calcário vira
um mármore por calor e pressão. Um granito se transforma em grãos de \omterm{def:g3:stones-pebbles-sand:sand}{areia}
por \omterm{def:g5:water-shapes-the-land:weathering}{intemperismo}.''
\end{solution}

\begin{solution}{exo:g6:rocks-and-minerals:5}
(a) Basalto. (b) Mármore. (c) Arenito.
\end{solution}

\begin{solution}{exo:g6:rocks-and-minerals:6}
A sua \omterm{def:g6:rocks-and-minerals:hardness}{dureza} é maior que a da unha (cerca de 2.5) e menor que a de uma
moeda de cobre (cerca de 3.5). Ele poderia ser a calcita (3).
\end{solution}

\begin{solution}{exo:g6:rocks-and-minerals:7}
Quartzo: $21 \div 60 = 0.35$, ou seja, 35~\%. Feldspato: $33 \div 60 = 0.55$,
ou seja, 55~\%. Mica: $6 \div 60 = 0.10$, ou seja, 10~\%. Verificação:
$35 + 55 + 10 = 100$.
\end{solution}

\begin{solution}{exo:g6:rocks-and-minerals:8}
O granito esfria muito devagar, no fundo da terra: os seus grãos têm muito
tempo para crescer, e ficam grandes. O basalto vem da \omterm{def:g4:volcanoes-earthquakes-first-look:eruption}{lava}, que esfria
depressa na superfície: os seus grãos não têm tempo de crescer e continuam
pequenos demais para ver.
\end{solution}

\begin{solution}{exo:g6:rocks-and-minerals:9}
(a) O calcário é sedimentar; o mármore é metamórfico. (b) No fundo da
terra, o calor e a pressão o transformaram sem derretê-lo: a sua calcita
cresceu em cristais brilhantes que se encaixam. (c) Os dois são feitos do
mesmo \omterm{def:g6:rocks-and-minerals:mineral}{mineral}, a calcita, e é a calcita que faz efervescência com ácido.
\end{solution}

\begin{solution}{exo:g6:rocks-and-minerals:10}
(a) Sim: o quartzo (7) é mais duro que o vidro (cerca de 5.5). (b) Não: a
fluorita (4) é mais mole que o vidro. (c) Sim: um prego de aço (cerca de
6.5) é mais duro que o feldspato (6).
\end{solution}

\begin{solution}{exo:g6:rocks-and-minerals:11}
Por exemplo: (1) o granito da \omterm{def:g1:landscapes-around-us:landform}{montanha} (ígneo) se fragmenta pelo
\omterm{def:g5:water-shapes-the-land:weathering}{intemperismo}, e os seus grãos de quartzo se soltam; (2) um \omterm{def:g1:landscapes-around-us:water}{rio} os leva
embora: a \omterm{def:g5:water-shapes-the-land:erosion}{erosão}; (3) os grãos são depositados em camadas no fundo do \omterm{def:g2:water-on-earth:ocean}{mar}:
a \omterm{def:g5:water-shapes-the-land:deposition}{deposição}; (4) enterrados embaixo de outras camadas, eles são prensados
e cimentados até virarem um arenito: o endurecimento (a \omterm{def:g3:stones-pebbles-sand:rock}{rocha} agora é
sedimentar); (5) levado para o fundo da terra, o arenito é transformado
pelo calor e pela pressão numa \omterm{def:g6:rocks-and-minerals:rock-families}{rocha metamórfica}.
\end{solution}

\begin{solution}{exo:g6:rocks-and-minerals:12}
(a) Não é riscado pelo aço (cerca de 6.5) e consegue riscar o vidro: mais
duro que 6.5. É riscado pelo topázio (8): mais mole que 8. A sua \omterm{def:g6:rocks-and-minerals:hardness}{dureza} é
7: é o quartzo. (b) O quartzo é mais duro que quase tudo o que se esfrega
nele (outros grãos, \omterm{def:g3:stones-pebbles-sand:pebble}{seixos}): ele quase não se desgasta. A calcita (3) e a
mica (2 a 3) são muito mais moles: o atrito as moi em pedacinhos que são
levados embora.
\end{solution}

\begin{solution}{pb:g6:rocks-and-minerals:1}
\textbf{1.} Granito, uma \omterm{def:g6:rocks-and-minerals:rock-families}{rocha ígnea}; os seus \omterm{def:g6:rocks-and-minerals:mineral}{minerais} são o feldspato
(rosado ou branco), o quartzo (cinza, vítreo) e a mica (preta, brilhante).

\textbf{2.} Calcário, uma \omterm{def:g6:rocks-and-minerals:rock-families}{rocha sedimentar}. As conchas pertenciam a animais
que viviam no \omterm{def:g2:water-on-earth:ocean}{mar}; os seus restos foram depositados e endurecidos na
\omterm{def:g3:stones-pebbles-sand:rock}{rocha}.

\textbf{3.} Ardósia, uma \omterm{def:g6:rocks-and-minerals:rock-families}{rocha metamórfica}, formada a partir do folhelho.

\textbf{4.} Mármore, uma \omterm{def:g6:rocks-and-minerals:rock-families}{rocha metamórfica}, formada a partir do calcário.

\textbf{5.} O quartzo: não é riscado pelo aço (cerca de 6.5) e risca o
vidro (cerca de 5.5); a sua \omterm{def:g6:rocks-and-minerals:hardness}{dureza} é 7.

\textbf{6.} Riscados pelo aço (cerca de 6.5), mas não pelo vidro (cerca de
5.5): entre 5.5 e 6.5. É o feldspato (ortoclásio, 6).

\textbf{7.} A mica (\omterm{def:g6:rocks-and-minerals:hardness}{dureza} de 2 a 3).

\textbf{8.} Sim: o calçamento é de calcário, feito de calcita (3), mais
mole que uma moeda de cobre (cerca de 3.5).

\textbf{9.} O \omterm{def:g5:water-shapes-the-land:weathering}{intemperismo}: a \omterm{def:g4:clouds-rain-snow:precipitation}{chuva} desgasta devagar a superfície da
pedra.

\textbf{10.} A \omterm{def:g1:landscapes-around-us:landform}{montanha} sofre \omterm{def:g5:water-shapes-the-land:weathering}{intemperismo}; os seus grãos de quartzo são
levados pelos \omterm{def:g1:landscapes-around-us:water}{rios} e depositados em camadas, no fundo de um \omterm{def:g1:landscapes-around-us:water}{lago} ou do \omterm{def:g2:water-on-earth:ocean}{mar};
enterrados e endurecidos, eles podem virar um arenito, uma \omterm{def:g6:rocks-and-minerals:rock-families}{rocha
sedimentar}.

\textbf{11.} As três famílias estão presentes: ígnea (o granito do
chafariz), sedimentar (o calcário do calçamento), metamórfica (a ardósia
do telhado e o mármore da estátua).

\textbf{12.} O \omterm{def:g6:rocks-and-minerals:mineral}{mineral} mais duro é o quartzo do granito: \omterm{def:g6:rocks-and-minerals:hardness}{dureza} de Mohs
\textbf{7}.
\end{solution}
